\documentclass[10pt,a4paper]{amsart}
\usepackage[margin=19mm]{geometry}
\usepackage{amsmath,amssymb,mathtools}
\usepackage[scr=rsfs,cal=euler]{mathalfa}
\usepackage{xfrac}
\usepackage[unicode,hidelinks,bookmarks=false]{hyperref}
\newcommand{\ie}{i.e.\ }
\newcommand{\cf}{cf.\ }
\newcommand{\resp}{respectively\ }
\newcommand{\Subsection}{\subsection}
\providecommand{\nobreakdash}{\nobreak\hbox{-}}
\setlength{\emergencystretch}{2em}
\multlinegap0pt
\usepackage{mathtools}
\usepackage[normalem]{ulem}
\usepackage{enumerate}
\usepackage[shortlabels]{enumitem}
\usepackage{cancel}
\usepackage{amssymb}
\usepackage{dsfont}
\usepackage{xcolor}
\newtheorem{lem}{Lemma}
\newtheorem{thm}{Theorem}
\title{Remarks on Lifespan and Continuation Criterion of Two Dimensional Incompressible Fluid Models}
\author{Anping Pan}
\date{Source: arXiv:2604.15721v2 (CC BY 4.0)}
\begin{document}
\maketitle

\noindent\emph{Source and licence.} Remarks on Lifespan and Continuation Criterion of Two Dimensional Incompressible Fluid Models. Version arXiv:2604.15721v2 (CC BY 4.0), distributed by arXiv under the Creative Commons Attribution 4.0 International licence, http://creativecommons.org/licenses/by/4.0/, as recorded in the arXiv licence metadata for this exact version and shown on the abstract page https://arxiv.org/abs/2604.15721v2. Full licence notice: \texttt{licenses/arxiv-2604.15721-CC-BY-4.0.txt}.\par\smallskip

\noindent\textit{Editorial scope.} Counted unit: the article's thm IIE (source0.tex lines 732--737). The article's complete original proof of it is delivered with it (source0.tex lines 738--856, one \texttt{proof} environment of 119 lines). No mathematics is written, repaired or re-derived: the statement and the proof are the article's own lines, in the article's own notation, and no retained line is substituted or re-flowed.  Retained uncounted as the source-local context the proof needs: lines 248--250; lines 340--351; lines 361--379; lines 416--442; lines 520--528; lines 621--626; lines 648--658; lines 681--683.  Nothing was omitted from any retained span, so the package carries no omission record.  No retained line contains a citation, so the wrapper carries no bibliography and no bibliography entry.  No retained line contains a figure environment, an image-inclusion command or drawing code, so no image is delivered and the package declares no asset dependency.  Editorial formatting only, in addition: an AMS \texttt{amsart} preamble that restates the article's own declarations and macros used by the retained lines, namely the environments \texttt{lem}, \texttt{thm} on separate counters, together with the standard packages those lines need.  The retained environments print in this document as Lemma 1, Theorem 1 in order of appearance, whereas the article prints the same items with the numbering of its own full document.  The complete native source of the article as distributed in the arXiv e-print for this exact version is bundled unchanged as \texttt{sources/198639/source0.tex}.
\begin{equation}\label{Streching Variable Definition}
M=\exp\bigg(C\int_0^t \lVert\nabla u\rVert_\infty d\tau\bigg),\quad N= \exp\bigg(C\int_0^t \lVert\nabla u\rVert_{1,p} d\tau\bigg),
\end{equation}

\begin{lem}\label{Bootstrap Lemma}
    Let $I$ be an interval of time. Suppose we have one hypothesis $\mathscr H(t)$ and one outcome $\mathscr O(t)$ defined for each $t\in I$, then $\mathscr O(t)$ holds true for all $t\in I$ provided the validity of the following statements:
    \begin{itemize}
        \item [(I)] Suppose $\mathscr H(t)$ is true at $t\in I$, then $\mathscr O(t)$ is true.

        \item [(II)] Suppose $\mathscr O(t)$ holds for $t\in I$, then there exists a neighborhood $I_t\owns t$ such that $\mathscr H(t^\prime)$ is true for all $t^\prime\in I_t$.

        \item [(III)] There exists a $t_*\in I$ such that $\mathscr H(t_*)$ is true.

        \item [(IV)] Assume $\{t_n\}_{n=1}^\infty$ is a sequence of time converging to $\tilde t$ such that $\mathscr O(t_j)$ is true for all $j$, then $\mathscr O(\tilde t)$ is true.
    \end{itemize}
\end{lem}

\begin{lem}\label{Hierarchical Growth of Norms for Bou and IIE}
  Let $M,N,Q: [0,T]\to \mathbb R_+$ be absolutely continuous non-decreasing functions, such that
  \begin{equation*}
  M(0)=1, \quad N(0)=1,\quad Q(0)=0.    
  \end{equation*}
  Suppose that there exists an interval $[0,T_*]$ and a constant $C>e$ such that the following holds on $[0,T_*]$:
  \begin{equation}\label{Differential inequality of MYZ}
   \dot M\le CM(1+\log(1+\frac{\dot N}{N})),\quad \dot N\le CN(1+M),  
  \end{equation}
  and
  \begin{equation}\label{Differential inequality of Duhamel Term}
    \dot Q\le C(\dot M+N\dot M/M)+M\dot M\dot N/N +C(M+N). 
  \end{equation}
  Then, there exists positive constants $\Lambda_1,\Lambda_2,\Lambda_3,\Lambda_4$ such that
  \begin{equation}\label{Growth of controlling norms}
  M(t)\le \exp(e^{\Lambda_1 t}),\quad N(t)\le \exp(\Lambda_2\exp(e^{\Lambda_1t})),\quad Q(t)\le \Lambda_4 \exp(\Lambda_3\exp( e^{\Lambda_1t})) .  
  \end{equation}
  
\end{lem}

\begin{lem}\label{Bootstrap Closure Lemma}
 Let $I=[0,T]$ and $\delta>0$, assume that $\{\mathcal F_j\}_{j=1}^m: I\to [0,+\infty)$ are continuous increasing functions with $\mathcal F_j(0)\le 1$ holding for all $1\le j\le m$, and we assume $\{\kappa_j\}_{j=1}^m$ is a sequence of positive real numbers. Let
 \begin{equation*}
  \mathcal E(t):=\exp(C_1\exp(C_2 e^{C_3t})),\quad C_1,C_2,C_3>0.   
 \end{equation*}
 Assume that there exists $C_1,C_2,C_3>0$ and sequence $\{\zeta_j\}_{j=1}^m \subset [1,+\infty)$ such that the following implication holds: For some $0<t\le T$:
 \begin{equation}\label{Bootstrap Implication}
\kappa_j\delta \mathcal F_j(t)\le 1,\quad\forall 1\le j\le m \quad\Longrightarrow\quad   \mathcal F_j(s)\le \zeta_j\mathcal E(s), \quad\forall 1\le j\le m, s\in [0,t].
 \end{equation}

 Then, for some small threshold 
 \begin{equation}\label{Delta0 Definition}
  \delta_0=C_0\exp(-C_1 e^{2C_2})  
 \end{equation}
 and $\delta<\delta_0$, we have: 
 \begin{equation*}
  \kappa_j\delta \mathcal F_j(t)\le 1  \quad\text{for all }1\le j\le m , t\in[0,T_\delta], 
 \end{equation*}
 where
 \begin{equation}\label{Triple Log Lifespan}
  T_\delta=C_3^{-1}\log[C_2^{-1}\log(C_1^{-1}\log(C_0\delta^{-1}))],\quad C_0:=\frac{99}{100\max_{1\le j\le m}  \kappa_j \zeta_j}.
 \end{equation}
 As a matter of fact, we have:
 \begin{equation}\label{Finiteness on Bootstrap interval}
   \mathcal F_j(t)  \le \zeta_j C_0\delta^{-1}\quad\text{for all }1\le j\le m, t\in[0,T_\delta].
 \end{equation}
\end{lem}

    \begin{equation}\label{Boussinesq Differential Inequalities}
\left\{
\begin{aligned}
&\dot M\le CM[1+\log(2+\frac{\dot N}{N})] (1+\delta Z),\\
&\dot N\le CN (1+M+M\delta Y),\\
&  \dot Y=M+N,\quad \dot Z=M,\quad Y_0=Z_0=0.
\end{aligned}
\right.
\end{equation}

\begin{lem}\label{Rho 0 Estimate}
Assume $\delta<1$, then we have
\begin{equation}
 \lVert\nabla\rho_0\rVert_\infty\lesssim \delta,\quad \lVert \nabla^2\rho_0\rVert_p\lesssim\delta.
\end{equation}
\end{lem}

\begin{lem}\label{Estimates of u}
Let $M(t),N(t)\ge 1$ be defined as in \eqref{Streching Variable Definition}, then
 there is a numerical constant $C\ge 1$ which satisfies the following: For any fixed  $\delta<C^{-1}$, there exists a corresponding time $T_\delta^*>0$ such that for all $t\in [0,T_\delta]$ the following estimates hold:
\begin{equation}\label{W2p estimate for u}
 \lVert u\rVert_{2,p}\le  \frac{ (1+\delta N)\lVert\omega\rVert_{1,p}}{1-C\delta N}
\end{equation}
\begin{equation}\label{Linfty estimate for nabla u}
 \lVert\nabla u\rVert_\infty\lesssim   \Big(\delta M\lVert\omega\rVert_\infty+\frac{ \lVert\omega\rVert_\infty}{1-C\delta M}\Big)\Big[1+\log 
 \Big(2+\frac{ (1+\delta N)\lVert\omega\rVert_{1,p}}{1-C\delta N}\Big)\Big].
\end{equation}
\end{lem}

\begin{equation}\label{u Wk,p Estimate}
 \le (1+\delta \lVert\theta\rVert_{k,p})\lVert\omega\rVert_{k-1,p}+(1+\delta \lVert\theta\rVert_{k,p})\frac{C_1\delta \lVert \theta\rVert_{k,p}\lVert \omega\rVert_{k-1,p}}{1-C_1\delta \lVert \theta\rVert_{k,p}}.   
\end{equation}

\begin{thm}\label{Lifespan of IIE}
There exist positive constants $\delta_0<1$ and $C_0^*,C_1^*, C_2^*,C_3^*$, possibly depending on the data $\omega_0$, such that for any $\delta<\delta_0$, $\mathsf E(t):=\lVert\omega\rVert_{1,p}+\lVert\rho\rVert_{2,p}$ remains bounded on the interval $[0,T_\delta]$, where
\begin{equation*}
 T_\delta=C_3^*\log[C_2^*\log( C_1^*\log (C_0^*\delta^{-1}))]. 
\end{equation*}
\end{thm}

\begin{proof}[Proof of Theorem~\ref{Lifespan of IIE}]
We divide the proof into several steps.
\begin{itemize}
    \item [(I)] We start with necessary estimates to establish the system of differential inequality, like \eqref{Boussinesq Differential Inequalities}. Consider the Duhamel term:
\begin{equation*}
\int_0^t g(\tau,X_\tau)d\tau=\nabla^\perp\rho_0 \int_0^t \nabla^* X_\tau (\nabla^* uu)(X_\tau)d\tau,
\end{equation*}
which admits the following estimates holding for generic $\delta\le 1$, thanks to Lemma \ref{Rho 0 Estimate}:
\begin{equation*}
 \bigg\lVert  \int_0^t g(\tau,X_\tau)d\tau \bigg\rVert_{1,p}\lesssim  \delta\int_0^\infty \lVert \nabla X_\tau\rVert_\infty \lVert\nabla u\rVert_\infty \lVert u\rVert_\infty d\tau+\delta  \int_0^t\lVert\nabla^2 X_\tau\rVert_p \lVert\nabla u\rVert_\infty \lVert u\rVert_\infty
\end{equation*}
\begin{equation*}
 +\delta\int_0^t\lVert\nabla X\rVert_\infty^2 \lVert\nabla^2 u\rVert_p\lVert \nabla u\rVert_\infty d\tau .  
\end{equation*}
Hence, we have:
\begin{equation*}
  \bigg\lVert  \int_0^t g(\tau,X_\tau)d\tau \circ A_t\bigg\rVert_{1,p} \lesssim  M\delta\Big[\int_0^t (\dot M+N\frac{\dot M}{M})\lVert u\rVert_\infty+\frac{d}{dt}M^2 \frac{\dot N}{N}d\tau
  \Big].
\end{equation*}
Therefore, Duhamel's formula yields the following estimate:
\begin{equation*}
 \lVert\omega\rVert_{1,p}\lesssim 1+M+  M\delta\Big[\int_0^t (\dot M+N\frac{\dot M}{M})\lVert u\rVert_\infty+M\dot M \frac{\dot N}{N}d\tau\Big],
\end{equation*}
and 
\begin{equation}\label{Energy norm of IIE}
\mathcal E(t)\lesssim  1+M +\delta M\Big[\int_0^t (\dot M+N\frac{\dot M}{M})\lVert u\rVert_\infty+M\dot M \frac{\dot N}{N}d\tau\Big]+\delta N.
\end{equation}
So, for any time $t\ge 0$ and any $\delta<1$, we have the $\mathbf{unconditional}$ estimate:
\begin{equation}\label{unconditional Duhamel estimate}
    \lVert\omega\rVert_{1,p}\lesssim 1+M+  \delta MQ,\quad Q=\int_0^t (\dot M+N\frac{\dot M}{M})\lVert u\rVert_\infty+\dot M \frac{\dot N}{N}d\tau .
\end{equation}
By Lemma \ref{Estimates of u}, we obtain that for $t\in [0,T_\delta]$ and $\delta<\delta_0\le (2C)^{-1}$:
\begin{equation*}
 \frac{\dot N}{N}\lesssim \frac{1+\delta N}{1-C\delta N} \bigg[1+M+  M\delta\Big(\int_0^t (\dot M+N\frac{\dot M}{M})\lVert u\rVert_\infty+M\dot M \frac{\dot N}{N}d\tau\Big) \bigg] . 
\end{equation*}
Moreover, taking $k=1$ in \eqref{u Wk,p Estimate}, we have
\begin{equation*}
\lVert u\rVert_\infty\lesssim \lVert u\rVert_{1,p}\lesssim \bigg(\frac{1+\delta M}{1-C\delta M}\bigg)\lVert\omega\rVert_p
\end{equation*}
\begin{equation*}
 \lesssim   \frac{1+\delta M}{1-C\delta M}\lVert\omega\rVert_\infty.
\end{equation*}
By noticing that the bounding constant in \eqref{unconditional Duhamel estimate} is some $C^\prime=C(\lVert\omega_0\rVert_{1,p})$, abusing the notation a little bit, we conclude that the following system of differential inequalities hold on $[0,T_\delta]$ for some $C$ depending on $\omega_0$ and the constant $C$ in Lemma \ref{Estimates of u}:
\begin{equation}\label{Differential Inequality System IIE}
\left\{
\begin{aligned}
&{\dot M}\le C  M\Big[\lVert\omega\rVert_\infty+\frac{\lVert\omega\rVert_\infty}{1-C\delta M}\Big]\Big(1+\log \Big(2+\frac{\dot N}{N}\Big)\Big);\\
&\dot N\le  CN\bigg(\frac{1+\delta N}{1-C\delta N}\bigg) \big(1+M+  \delta MQ \big);  \\
&\dot Q=  \Big(\dot M+N\frac{\dot M}{M}\Big)\lVert u\rVert_\infty+M\dot M \frac{\dot N}{N},\quad Q(0)=0 . 
\end{aligned}
\right.    
\end{equation}
Moreover, as long as $C\delta M(t)<1$, we have the following estimate for $\lVert\omega\rVert_\infty$:
\begin{equation*}
 \lVert\omega(t,\cdot)\rVert_\infty \le \lVert\omega_0\rVert_\infty+\lVert\nabla \rho_0\rVert_\infty\int_0^t \lVert\nabla^*X\rVert_\infty \lVert\nabla u\rVert_\infty \lVert u\rVert_\infty d\tau  
\end{equation*}
\begin{equation}\label{Linfty estimate for omega}
 \le   C+C\delta \int_0^t \frac{(1+\delta M)\dot M}{1-C\delta M}\lVert\omega\rVert_\infty d\tau.
\end{equation}

\item [(II)]

Now we propose our bootstrap scheme. We say $\mathscr H(t)$ is satisfied if
\begin{equation}\label{Hypothesis IIE}
 C\delta M\le 1/2,\quad C\delta N\le 1/2  ,\quad \delta Q\le 1.
\end{equation}
Assuming that $\mathscr H(T_*)$ holds, then clearly $T_*\le T_\delta^*$ and the estimates \eqref{W2p estimate for u} and \eqref{Linfty estimate for nabla u} hold. Since $M,N,Q$ are all increasing, we conclude that $\mathscr H(t)$ holds for all $0\le t\le T_*$. 

Moreover, for $t\in [0,T_*]$, \eqref{Linfty estimate for omega} implies that
\begin{equation*}
\lVert\omega(t,\cdot)\rVert_\infty\le   C+3C\delta \int_0^t {\dot M(\tau)\lVert\omega(\tau,\cdot)\rVert_\infty} d\tau.  
\end{equation*}
By the integral version of Gr\"onwall inequality, we have:
\begin{equation*}
 \lVert\omega(t,\cdot)\rVert_\infty\le  C+3C^2\delta\int_0^t \dot M(\tau)\exp[3C\delta(M(t)-M(\tau))]d\tau
\end{equation*}
\begin{equation*}
 \le C+2C\exp(3C\delta M(t))\le 3e^2 C .   
\end{equation*}

Using the fact that $C\ge 1$, we are able to reduce \eqref{Differential Inequality System IIE} $[0,T_*]$ to the following system:
\begin{equation}\label{Relaxed Differential Inequality System IIE}
\left\{
\begin{aligned}
&{\dot M}\le  9e^2CM\Big(1+\log\Big(1+\frac{\dot N}{N}\Big)\Big);\\
&\dot N\le 6CN\big(1+M);  \\
&\dot Q\le 9e^2C(\dot M+N\dot M/M)+M\dot M\dot N/N.
\end{aligned}
\right.    
\end{equation}

We now introduce the following outcome $\mathscr O(t)$: We say $\mathscr O(t)$ holds if 
 there exists $\Lambda_1,\Lambda_2,\Lambda_3,\Lambda_4>0$ such that 
\begin{equation*}
M(t)\le \exp(e^{\Lambda_1 t}),\quad N(t)\le \exp(\Lambda_2\exp(e^{\Lambda_1t})),\quad Q(t)\le \Lambda_4\exp(\Lambda_3\exp(e^{\Lambda_1 t}).
\end{equation*}
Apply Lemma \ref{Hierarchical Growth of Norms for Bou and IIE} to \eqref{Relaxed Differential Inequality System IIE}, with the same $M$,$N$,$Q$ and constant $9e^2C$, it follows that $\mathscr H(t)$ implies $\mathscr O(t)$. Since $M,N,Q$ are all absolutely continuous and increasing with $M(0)=N(0)=1$, $Q(0)=0$, there exists sufficiently small $t$ such that $\mathscr H(t)$ holds. On the other hand, the closedness condition trivially holds for $\mathscr O(t)$.

\item [(III)] We now close the bootstrap, which amounts to show (II) in lemma \ref{Bootstrap Lemma}. To this end, apply Lemma \ref{Bootstrap Closure Lemma} with
\begin{equation*}
  m=3,\quad\mathcal F_1=M,\quad \mathcal F_2=N,\quad \mathcal F_3=Q,\quad \kappa_1=\kappa_2=1/2,\quad \kappa_3=1.  
\end{equation*}
and
\begin{equation*}
 \zeta_1=\zeta_2=1,\quad \zeta_3=   \Lambda_4.
\end{equation*}
As a consequence, we have the bootstrap closed on $[0,T]$ with $T\ge T_\delta$. More precisely:
\begin{equation*}
C_0=\frac{99}{100\max\{1/2,\Lambda_4\}},\quad C_1=\max\{\Lambda_2,\Lambda_3\},\quad C_2=1,\quad C_3=\Lambda_1    
\end{equation*}
and $\delta_0$, $T_\delta$ respectively defined by \eqref{Delta0 Definition} and \eqref{Triple Log Lifespan}. As a consequence, we have for all $t\in [0,T_\delta]$:
\begin{equation*}
  \mathsf E(t)\lesssim 1+M+\delta MQ +\delta N\lesssim 1+M\lesssim 1+C_0\delta^{-1}. 
\end{equation*}
Which completes the proof of the theorem.

\end{itemize}

\end{proof}

\end{document}
